\chapter{Kategorifeilen}

\begin{motto}
Eit fag svarar ikkje for retten til å eksistere\\ ved å bli eit anna fag.
\end{motto}

Heile saka mot designfaget kviler på éin premiss, og dette kapitlet handlar om den. Premissen er at eit fag berre har eit fundament dersom det kan grunngje dommane sine slik naturvitskapen gjer — med presise, falsifiserbare påstandar, kumulative resultat, ein delt og overførbar metode. Måler ein design mot den standarden, kjem faget til kort, og konklusjonen følgjer: design manglar eit fundament. Men konklusjonen følgjer berre dersom premissen held, og premissen held ikkje. Å krevje at design skal grunngje seg sjølv slik fysikken gjer, er ikkje eit krav om strenghet. Det er ein kategorifeil — ein feil av same slag som å klage på at eit dikt ikkje let seg etterprøve, eller at ein rettsavgjerd ikkje kan replikerast i eit laboratorium.

\section{Skiljet Cross gjorde}

Det var Nigel Cross som gav dette den klaraste forma, i ein argumentasjon som strekkjer seg frå den korte programartikkelen «Designerly Ways of Knowing» \autocite{cross1982} til den seinare utgreiinga \autocite{cross2006}. Cross sitt poeng er at det finst tre breie kulturar for kunnskap, ikkje éin. Det er vitskapen, hvis objekt er den naturlege verda og hvis metode er det kontrollerte eksperimentet. Det er humaniora, hvis objekt er den menneskelege erfaringa og hvis metode er tolkinga. Og det er design — eller det «kunstige», det menneskeskapte — hvis objekt er det som \emph{kunne vore annleis}, og hvis metode er ein eigen: å gje form til det som enno ikkje finst. Kvar kultur har sine eigne kriterium for kva som tel som dugleik. Å måle den eine med den andre sine kriterium er ikkje strenghet; det er forveksling.

Det avgjerande i Cross sitt skifte — frå draumen om ein \emph{designvitskap} til prosjektet om ein \emph{designdisiplin} — er at det ikkje er ei innrømming av nederlag. Det er ei korrigering av spørsmålet. Simon hadde drøymt om å gjere design til ein vitskap om det kunstige, modellert på naturvitskapen \autocite{simon1969sciences}; Cross sitt poeng er at dette var å be design om å bli noko anna enn det er. Disiplinen design har sin eigen epistemologi, sin eigen praksiologi, sin eigen fenomenologi — kunnskap om designarar, om designprosessar, og om designa ting — og desse let seg studere på sine eigne premissar. At dei ikkje liknar fysikken, er ikkje ein mangel ved design. Det er det som gjer design til ein eigen kultur.

\section{Kva ein kategorifeil er}

Lat oss vere presise om kva skuldinga inneber, for ho er lett å misforstå. Ein kategorifeil er ikkje berre ein usemje om standardar. Det er å bruke eit omgrep i ein samanheng der det ikkje gjev meining — å spørje kva farge talet sju har, eller kor tung rettferda er. Når kritikaren spør kvifor designforskinga ikkje akkumulerer slik biologien gjer, og tek fråvêret av biologisk-type akkumulasjon som prov på mangel, gjer han ein slik feil. Han har alt avgjort, i sjølve spørsmålet, at den einaste forma for kunnskap er den naturvitskaplege; og så «oppdagar» han at design ikkje har den forma. Men det var ikkje ei oppdaging. Det var ein føresetnad, forkledd som ein konklusjon.

Tenk på kor mange respekterte felt som ville falle for same testen. Jussen kan ikkje grunngje dommane sine slik fysikken kan; to dyktige juristar kan komme til motsette konklusjonar av same lov, kvar med fullgod grunngjeving, og likevel tvilar ingen på at jussen er ein disiplin med eit fundament. Klinisk medisin, i møtet med den einskilde pasienten, er full av skjøn som ikkje let seg redusere til protokoll. Historiefaget akkumulerer ikkje mot éi endeleg sanning. Arkitekturen, pedagogikken, psykiatrien — alle desse er modne fag med skjøn i kjernen, og ingen av dei oppfyller den standarden design vert målt mot. Anten må kritikaren felle alle desse faga med same slaget, eller så må han innrømme at det finst meir enn éi form for fagleg fundament. Han kan ikkje ha det begge vegar.

\section{Den uangripelege kritikken}

Det finst ein djupare ironi her, og han er verdt å peike på, fordi han snur kritikken mot seg sjølv. \textsc{Grunnlaust} hevdar at design ikkje kan falsifiserast, ikkje kan ta feil, og at dette er faget sin store veikskap. Men sjå på sjølve tesen om at design manglar eit fundament. Kva ville telje som å motbevise henne? Kvar gong forsvararen peikar på ein metode, vert det svart at metoden ikkje er ekte metode. Kvar gong forsvararen peikar på kumulativ kunnskap — mønster, prinsipp, retningslinjer — vert det svart at det ikkje er ekte kunnskap. Kvar gong forsvararen peikar på ein etikk-kodeks, vert det svart at han manglar handheving og difor ikkje tel. Tesen er konstruert slik at ingenting kan telje imot henne: alt som finst, vert omdefinert til ikkje å vere det rette slaget. Det er den same uangripelegheita kritikaren skuldar design for. Eit syn som immuniserer seg mot kvart moteksempel ved å omdefinere kva som ville telje, er ikkje ein streng vitskap. Det er ein sirkel.

Dette gjer ikkje at kritikken er verdlaus — han peikar på reelle tilhøve, som dei følgjande kapitla vil ta på alvor. Men det flyttar bevisbyrda. Kritikaren må vise, utan å føresetje det han skal prove, at den naturvitskaplege forma for fundament er den \emph{einaste} forma eit fag kan ha. Og det kan han ikkje, for verda er full av modne fag som motbeviser det. Resten av denne boka viser kva for eit fundament design faktisk har — ikkje fysikken sitt, men sitt eige.

\vignett

Om kravet er feilstilt, må vi spørje på nytt kva eit fag eigentleg er, og kva det vil seie å ha eit fundament. Det er det neste kapitlet.

\vidarelesing{\fullcite{cross2006}; \fullcite{cross1982}; \fullcite{simon1969sciences}; \fullcite{schon1983}; \fullcite{kuhn1962}.}
